%----------------------------------------------------------------------------------------
%    PACKAGES AND THEMES
%----------------------------------------------------------------------------------------

\documentclass[aspectratio=169,xcolor=dvipsnames]{beamer}

\AtBeginSection[]
{
    \begin{frame}
        \frametitle{Visão Geral}
        \tableofcontents[currentsection]
    \end{frame}
}

\usefonttheme{professionalfonts} % using non standard fonts for beamer
\usefonttheme{serif} % default family is serif
\usepackage{fontspec}
\setmainfont{XCharter}[
    UprightFont    = *-Roman,
    BoldFont       = *-Bold,
    ItalicFont     = *-Italic,
    BoldItalicFont = *-BoldItalic,
]

\usepackage{hyperref}
\usepackage{graphicx} % Allows including images
\usepackage{booktabs} % Allows the use of \toprule, \midrule and \bottomrule in tables
\input{./slides/SimplePlusTheme/beamerthemeSimplePlus.sty}
\input{./slides/SimplePlusTheme/beamerinnerthemeSimplePlus.sty}
\input{./slides/SimplePlusTheme/beamerfontthemeSimplePlus.sty}
\input{./slides/SimplePlusTheme/beamercolorthemeSimplePlus.sty}

\usepackage{listings}
\setbeamercovered{transparent}

\usepackage{biblatex}
\addbibresource{slides/reference.bib}
\renewcommand*{\bibfont}{\footnotesize}
% Tamanho dos slides de "Leituras Recomendadas": maior que a bibliografia
% completa (\bibfont acima, que rende \tiny via a redefinição de \footnotesize)
% para destacar as leituras curadas.
\newcommand{\leiturasize}{\small}

\usepackage[brazilian ]{babel}
\usepackage{csquotes}
\usepackage{pgfplots}
\pgfplotsset{compat=1.18}

\usepackage{emoji}
\usepackage{amsmath}
\usepackage[xcharter,bigdelims,vvarbb]{newtxmath}
% newtxmath's operators font requests the fontspec family `XCharter(0)' in OT1;
% without these substitutions math digits and operator names silently fall back
% to Computer Modern (LaTeX Font Warning: Font shape `OT1/XCharter(0)/m/n' undefined).
\DeclareFontFamily{OT1}{XCharter(0)}{}
\DeclareFontShape{OT1}{XCharter(0)}{m}{n}{<->ssub * XCharter-TLF/m/n}{}
\DeclareFontShape{OT1}{XCharter(0)}{m}{it}{<->ssub * XCharter-TLF/m/it}{}
\DeclareFontShape{OT1}{XCharter(0)}{b}{n}{<->ssub * XCharter-TLF/b/n}{}
\DeclareFontShape{OT1}{XCharter(0)}{bx}{n}{<->ssub * XCharter-TLF/b/n}{}

\usepackage{bm}
\usepackage[table]{xcolor}
\usepackage{xcolor}
\usepackage{algpseudocode}

\usepackage{cancel}
\usepackage{ulem}

\usetikzlibrary{shapes}
\usetikzlibrary{arrows}
\usetikzlibrary {arrows.meta}
\usetikzlibrary{calc,positioning}
\usepgfplotslibrary{fillbetween}
\usetikzlibrary{angles,quotes}
\usetikzlibrary{tikzmark}


\definecolor{PastelRed}{HTML}{FFC1CC}
\definecolor{PastelBlue}{HTML}{C2F0FF}
\definecolor{PastelBlue2}{HTML}{3182BD}

\definecolor{PastelPink}{HTML}{FFCBE1}
\definecolor{PastelGreen}{HTML}{D6E5BD}
\definecolor{PastelYellow}{HTML}{F9E1A8}
\definecolor{PastelBlue}{HTML}{BCD8EC}
\definecolor{PastelPurple}{HTML}{DCCCEC}
\definecolor{PastelOrange}{HTML}{FFDAB4}


\renewcommand{\footnotesize}{\tiny}

\newcommand{\mespor}{
  \ifcase\month
    \or janeiro
    \or fevereiro
    \or março
    \or abril
    \or maio
    \or junho
    \or julho
    \or agosto
    \or setembro
    \or outubro
    \or novembro
    \or dezembro
  \fi
}
% \usepackage{csquotes}
% \usepackage{minted}
% \usemintedstyle{xcode}
\definecolor{vscLightBackground}{HTML}{FFFFFF}
\definecolor{vscLightText}{HTML}{000000}
\definecolor{vscLightGreen}{HTML}{008000}        % For Comments AND Numbers
\definecolor{vscLightRed}{HTML}{A31515}          % For Strings
\definecolor{vscLightTeal}{HTML}{267F99}
\definecolor{vscBlue}{HTML}{0000FF}% For main keywords (if, for, def...)
\definecolor{vscLightPytorch}{HTML}{800080}      % For PyTorch keywords (dark magenta)

\lstdefinestyle{vscode-light}{
    language=Python,
    backgroundcolor=\color{vscLightBackground},
    basicstyle=\ttfamily\scriptsize\color{vscLightText},
    keywordstyle=\color{vscLightTeal},
    commentstyle=\color{vscLightGreen},
    stringstyle=\color{vscLightRed},
    % Style for PyTorch keywords (using keyword group 2)
    keywordstyle=[2]\color{vscLightPytorch}, 
    morekeywords={self, True, False, None},
    morekeywords=[2]{
        torch, Tensor, device, dtype, to, nn, optim, autograd, utils, data,
        Module, Linear, Conv2d, ReLU, CrossEntropyLoss, Adam, SGD,
        Dataset, DataLoader, randn, zeros, ones, arange, cat, stack,
        view, reshape, sum, mean, max, min, matmul, no_grad, backward
    },
    keywordstyle=[3]\color{vscBlue}, 
    morekeywords=[3]{
        model, loss, y, y_hat, x , optimizer
    },
    frame=single,
    tabsize=1,
    showstringspaces=false,
    breaklines=true,
    captionpos=b,
    extendedchars=true,
    numbers=left,
}

\usepackage[cache=true]{minted}
% minted v3+ (TeX Live 2024+) auto-detects the output directory.
% minted v2 (Ubuntu apt TeX Live 2023) needs it set explicitly to
% match latexmkrc's $out_dir = 'build'.
\makeatletter
\@ifpackagelater{minted}{2024/01/01}{}{%
  \def\minted@outputdir{build/}%
}
\makeatother
\setminted{
    fontsize=\scriptsize,
    style=vs,
    breaklines=true,
    numbers=left,
    numbersep=5pt,
    frame=single,
}
\providecommand{\citeauthoryear}[1]{(\citeauthor{#1},~\citeyear{#1})}

\providecommand{\thetav}{\bm{\theta}}

% vetor coluna 2D com colchetes
\providecommand{\cvec}[2]{\begin{bmatrix} #1 \\ #2 \end{bmatrix}}

\providecommand{\varvector}[1]{\mathbf{#1}}
% instance
\providecommand{\X}{\varvector{X}}
\providecommand{\mlinstance}{\varvector{X}}
\providecommand{\mlinstancei}{\varvector{X}^{(i)}}
\providecommand{\mlinstancex}[1]{\varvector{X}^{(#1)}}

% target
\providecommand{\y}{\varvector{y}}
\providecommand{\mltarget}{\varvector{y}}
\providecommand{\mltargeti}{\varvector{y}^{(i)}}
\providecommand{\mltargetx}[1]{\varvector{y}^{(#1)}}

%hats
\providecommand{\yhat}{\hat{\varvector{y}}}
\providecommand{\mltargethat}{\hat{\varvector{y}}}
\providecommand{\mltargethati}{\hat{\varvector{y}}^{(i)}}
\providecommand{\mltargethatx}[1]{\hat{\varvector{y}}^{(#1)}}

% linear reg
\providecommand{\mllinearregression}{\mltargethat=\mlinstance\bm{\theta}}

\providecommand{\tikzarrowcolor}[3]{%
  \begin{tikzpicture}[remember picture, overlay]
    % Arrow body
    \path[fill=#3] ([shift={(#1,#2)}]current page.south west) 
      rectangle ++(0.3, -0.2);
    % Arrow head
    \path[fill=#3] ([shift={(#1+0.3,#2+0.1)}]current page.south west) -- 
         ++(0, -0.4) -- 
         ++(0.2, 0.2) -- 
         cycle;
  \end{tikzpicture}%
}

\providecommand{\tikzarrow}[2]{%
    \tikzarrowcolor{#1}{#2}{Red}
}



\DeclareMathOperator*{\argmax}{argmax}
\DeclareMathOperator*{\argmin}{argmin}



%----------------------------------------------------------------------------------------
%    TITLE PAGE
%----------------------------------------------------------------------------------------

\title{Deep Learning}
\subtitle{Neurônio MP e Perceptron}

\author{Lucas Silveira Kupssinskü}

\institute
{
    Machine Learning Theory and Applications Laboratory \\
    PUCRS % Your institution for the title page
}
\date{\mespor de \the\year}
 % Date, can be changed to a custom date

%----------------------------------------------------------------------------------------
%    PRESENTATION SLIDES
%----------------------------------------------------------------------------------------

\begin{document}

\begin{frame}
    % Print the title page as the first slide
    \titlepage
\end{frame}

\begin{frame}{Overview}
    % Throughout your presentation, if you choose to use \section{} and \subsection{} commands, these will automatically be printed on this slide as an overview of your presentation
    \tableofcontents
\end{frame}
\section{Introdução}
%------------------------------------------------
\begin{frame}{Bioinspiração}
\begin{columns}
    \begin{column}{.5\linewidth}
        \begin{itemize}
            \item Quando não sabemos como fazer algo podemos copiar!
            \item Temos diversos exemplos de inteligência na natureza
            \item Computação Bioinspirada
            \begin{itemize}
                \item Algoritmos Genéticos 
                \item Computação Evolutiva
                \item Redes Neurais Artificiais...
            \end{itemize}
        \end{itemize}
    \end{column}
    \begin{column}{.5\linewidth}
          \centering
          \includegraphics[width=\linewidth]{slides/img/aula10/pessoavoando.jpeg}
    \end{column}
\end{columns}

\end{frame}

\section{Neurônio MP}

\begin{frame}{MP Neuron\footfullcite{mcculloch1943logical}}
\begin{columns}

    \begin{column}{.7\linewidth}
        \begin{itemize}
            \item Importância histórica por ser o primeiro modelo computacional de um neurônio
            \item Possibilitou mapear a lógica proposicional em neurônios artificiais
            \item Possível origem do termo \textit{neurônio artificial} e \textit{redes neurais artificiais}
            \item Modelo matemático, sem implementação concreta em um dispositivo computacional
        \end{itemize}
    \end{column}
    \begin{column}{.3\linewidth}
    \begin{figure}
        \centering
        \includegraphics[width=\linewidth]{slides/img/aula10/mculochpits.png}
        \caption{Warren S. McCulloch (à direita) e Walter Pitts (à esquerda).}
    \end{figure}
    \end{column}
\end{columns}
\end{frame}

\begin{frame}{MP Neuron – Funcionamento}
\begin{columns}
    \begin{column}{.6\textwidth}
        \begin{itemize}
            \item Entradas binárias
            \item Estímulos podem ser:
            \begin{itemize}
                \item Excitatórios: Ativos quando a entrada é Verdadeira
                \item Inibitórios: Ativos quando a entrada é Falsa
            \end{itemize}
            \item Neurônio dispara quando a quantidade de estímulos excitatórios ultrapassa um limiar $\theta$
        \end{itemize}
    \end{column}
    \begin{column}{.4\linewidth}
    \resizebox{\linewidth}{!}{
        \providecommand{\figxzerovalue}{$\mlinstance_{0}$}
        \providecommand{\figxonevalue}{$\mlinstance_{1}$}
        \providecommand{\figyhatvalue}{$\mltargethat$}
        \input{slides/img/tikz/mpneuron.tikz}
        }
    \end{column}
\end{columns}
\end{frame}

\begin{frame}{MP Neuron – Exemplo como \textit{AND}}
    \centering
    \resizebox{.6\linewidth}{!}{
        \providecommand{\figxzerovalue}{$\mlinstance_{0}=\textcolor{blue}{\text{true}}$}
        \providecommand{\figxonevalue}{$\mlinstance_{1}=\textcolor{red}{\text{false}}$}
        \providecommand{\figyhatvalue}{$\mltargethat=\textcolor{red}{\text{false}}$}
        \providecommand{\figpreactivation}{$2$}
        \input{slides/img/tikz/mpneuron.tikz}
        }
\end{frame}

\begin{frame}{MP Neuron – Exemplo como \textit{OR}}
    \centering
    \resizebox{.6\linewidth}{!}{
        \providecommand{\figxzerovalue}{$\mlinstance_{0}=\textcolor{blue}{\text{true}}$}
        \providecommand{\figxonevalue}{$\mlinstance_{1}=\textcolor{red}{\text{false}}$}
        \providecommand{\figyhatvalue}{$\mltargethat=\textcolor{blue}{\text{true}}$}
        \providecommand{\figpreactivation}{$1$}
        \input{slides/img/tikz/mpneuron.tikz}
        }
\end{frame}

\begin{frame}{MP Neuron – Entrada inibitória $\mlinstance_{0}\  \text{\textit{AND}}\ \text{NOT}\ \mlinstance_{1}$}
    \centering
    \resizebox{.6\linewidth}{!}{
        \providecommand{\figxzerovalue}{$\mlinstance_{0}=\textcolor{blue}{\text{true}}$}
        \providecommand{\figxonevalue}{$\mlinstance_{1}=\textcolor{red}{\text{false}}$}
        \providecommand{\figyhatvalue}{$\mltargethat=\textcolor{blue}{\text{true}}$}
        \providecommand{\figpreactivation}{$2$}
        \begin{tikzpicture}
    \providecommand{\figpreactivation}{$\theta$}
    \providecommand{\figxzerovalue}{$\mlinstance_{0}=\text{true}$}
    \providecommand{\figxonevalue}{$\mlinstance_{1}=\text{true}$}
    \providecommand{\figyhatvalue}{$\mltargethat=\text{true}$}
    \pgfmathsetmacro{\ex}{0.5}
    \pgfmathsetmacro{\ey}{1}
    \pgfmathsetmacro{\eradius}{1}
    \pgfmathsetmacro{\yone}{1.5}
    \pgfmathsetmacro{\ythree}{-1.5}
    % \draw[help lines] (-2,-2) grid (3,3);
    % first
    \draw[thick] (-2, \yone) -- (\ex, \ey-1);
    \fill (-2, \yone) circle [x radius=0.15, y radius=0.15];    
    \node at (-2, \yone+0.5) {\large \figxzerovalue};

    % third
    \draw[thick] (-2, \ythree) -- (\ex, \ey-1);
    \fill (-2, \ythree) circle [x radius=0.15, y radius=0.15]; 
    \filldraw[fill=white] (-0.5, \ythree+0.9) circle [x radius=0.15, y radius=0.15]; 
    
    \node at (-2, \ythree-0.3) {\large \figxonevalue};

    
    \draw[fill=PastelBlue] (\ex,\ey) arc(90:270:\eradius) --cycle;
    \draw[fill=PastelBlue2] (\ex,\ey) arc(90:-90:\eradius) --cycle;
    
    \node[font=\Large] at (\ex - 0.4, 0) {\huge \figpreactivation}; %

    \draw[-{Latex[length=3mm]},thick] (\ex+1, 0) -- (\ex+2, 0);
    \node at (\ex+2, 0.3) {\large \figyhatvalue};
\end{tikzpicture}
        }
\end{frame}

\section{Perceptron}

\begin{frame}{Perceptron}

Frank Rosenblatt\footfullcite{rosenblatt1958perceptron} fez algumas modificações no Neurônio MP
\begin{equation*}
    \mltargethat = \varphi(\mlinstance\thetav^{\top})
\end{equation*}
\begin{itemize}
    \item Entradas são números reais
    \item Cada entrada é ponderada por um parâmetro
    \item Existe um algoritmo de aprendizado para definir $\thetav$
    \item Função de ativação $\varphi(z)=\text{sinal}(z)=\begin{cases}
        +1,\quad z\geq0\\
        -1,\quad z<0
    \end{cases}$ 
\end{itemize}
\end{frame}

\begin{frame}{Perceptron}

\large ``(...) identities of this sort must be \textit{learned} (...) the number of functional units in the storage system or memory, should be less than the number of forms or memories to be retained.''\footfullcite{rosenblatt1958perceptron}
\end{frame}

\begin{frame}{Perceptron}
\begin{columns}
    \begin{column}{.6\textwidth}
        \begin{itemize}
            \item Classificador de imagens
            \item Resolução $20\times20$
            \item 400 fotocélulas capturavam a imagem
            \item Pesos eram controlados por potenciômetros 
        \end{itemize}
        
    \end{column}
    \begin{column}{.4\textwidth}
        \begin{figure}
            \centering
            \includegraphics[width=\linewidth]{slides/img/aula10/perceptron_rosenblatt.jpg}
            \caption{\href{https://news.cornell.edu/stories/2019/09/professors-perceptron-paved-way-ai-60-years-too-soon}{Frank Rosenblatt e o Perceptron}{}}
        \end{figure}
    \end{column}
\end{columns}

\end{frame}

\begin{frame}{Algoritmo de Aprendizado do Perceptron (PLA)}
\begin{algorithmic}[1]
\State Assuma $\X^{(i)}$ com primeira componente igual a $1$ (truque do bias)
\State Inicialize os pesos $\thetav$ aleatoriamente
\While{existirem observações mal classificadas}
  \State Selecione $\X^{(i)}$ tal que $\text{sinal}(\X^{(i)} \thetav^\top) \neq \y^{(i)}$
  \State Atualize: $\thetav \gets \thetav + \y^{(i)} \X^{(i)}$
\EndWhile
\end{algorithmic}
\end{frame}

\begin{frame}{Algoritmo de Aprendizado do Perceptron (PLA)}
\begin{columns}
    \begin{column}{.6\textwidth}
    \begin{itemize}
        \item Vamos ver um exemplo de treinamento com base no conjunto de dados abaixo
        \item À direita teremos uma representação do espaço de \textit{features} junto da fronteira de decisão
    \end{itemize}
        \begin{table}
            \centering
            \begin{tabular}{|r|r|r|}
                \hline
                $\X_1$ & $\X_2$  & $\y$  \\
                \hline
                1 & 1&\textcolor{blue}{1} \\
                \hline
                -1 & 1&\textcolor{blue}{1} \\
                \hline
                1 & -1&\textcolor{blue}{1} \\
                \hline
                -1 & -1&\textcolor{red}{-1} \\
                \hline
            \end{tabular}
        \end{table} 
    \end{column}
    \begin{column}{.4\textwidth}
    \centering
    \providecommand{\figfuncdecision}{x*x}
    \resizebox{\linewidth}{!}{
        \begin{tikzpicture}
\providecommand{\figfuncdecision}{-x+1}
\providecommand{\figtopcolor}{PastelRed}
\providecommand{\figbottomcolor}{PastelBlue}
\providecommand{\figvertical}{}
\providecommand{\figleftcolor}{PastelBlue}
\providecommand{\figrightcolor}{PastelRed}
\def\figemptytest{}
\begin{axis}[
    legend pos=south east,% <-- this line controls the width
    xlabel={$\mlinstance_1$},
    ylabel={$\mlinstance_2$},
    xmin=-1.5, xmax=1.5,
    ymin=-1.5, ymax=1.5,
    xtick={-1, 0, 1},
    ytick={-1, 0, 1},
    ymajorgrids=true,
    xmajorgrids=true,
    grid style=dashed,
]
    \ifx\figvertical\figemptytest
    \addplot [
        name path=f,
        domain=-1.5:1.5,
        samples=2,
        color=black,
        dashed,
        thick,
        ]
        {\figfuncdecision};
    \path [name path=bottom](\pgfkeysvalueof{/pgfplots/xmin},\pgfkeysvalueof{/pgfplots/ymin}) --
    (\pgfkeysvalueof{/pgfplots/xmax},\pgfkeysvalueof{/pgfplots/ymin});

    \path [name path=top]
    (\pgfkeysvalueof{/pgfplots/xmin},\pgfkeysvalueof{/pgfplots/ymax}) --
    (\pgfkeysvalueof{/pgfplots/xmax},\pgfkeysvalueof{/pgfplots/ymax});


    \addplot[\figtopcolor, opacity=0.3] fill between [of=f and top];
    \addplot[\figbottomcolor, opacity=0.3] fill between [of=f and bottom];
    \else
    \fill[\figleftcolor, opacity=0.3] (axis cs:-1.5,-1.5) rectangle (axis cs:\figvertical,1.5);
    \fill[\figrightcolor, opacity=0.3] (axis cs:\figvertical,-1.5) rectangle (axis cs:1.5,1.5);
    \draw[black, dashed, thick] (axis cs:\figvertical,-1.5) -- (axis cs:\figvertical,1.5);
    \fi

    \addplot[
        only marks,
        ultra thick,
         color=red,
         mark=-,
         mark options={scale=3}
        ]
        coordinates {
        (-1,-1)
        };
    \addplot[
        only marks,
        ultra thick,
         color=blue,
         mark=+,
         mark options={scale=3}
        ]
        coordinates {
        (1,1)(-1,1)(1,-1)
        };
\end{axis}
\end{tikzpicture}

    }
    \end{column}
\end{columns}
\end{frame}

\begin{frame}{Algoritmo de Aprendizado do Perceptron (PLA)}

\begin{columns}
    \begin{column}{.6\textwidth}
        \begin{algorithmic}[1]
            \State Assuma $\X^{(i)}$ com primeira componente igual a $1$ (truque do bias)
            \State Inicialize os pesos $\thetav$ aleatoriamente
            \While{existirem observações mal classificadas}
              \State Selecione $\X^{(i)}$ tal que $\text{sinal}(\X^{(i)} \thetav^\top) \neq \y^{(i)}$
              \State Atualize: $\thetav \gets \thetav + \y^{(i)} \X^{(i)}$
            \EndWhile
        \end{algorithmic}
    \centering
    \resizebox{.5\linewidth}{!}{
        \begin{tikzpicture}
    \providecommand{\figpreactivation}{$z^{(i)}$}
    \providecommand{\figactivation}{$a^{(i)}$}
    \providecommand{\figthetazero}{$\bm{\theta_{0}}$}
    \providecommand{\figthetaone}{$\bm{\theta_{1}}$}
    \providecommand{\figthetatwo}{$\bm{\theta_{2}}$}
    \pgfmathsetmacro{\ex}{0.5}
    \pgfmathsetmacro{\ey}{1}
    \pgfmathsetmacro{\eradius}{1}
    \pgfmathsetmacro{\yone}{1.5}
    \pgfmathsetmacro{\ytwo}{0}
    \pgfmathsetmacro{\ythree}{-1.5}
    % \draw[help lines] (-2,-2) grid (3,3);
    % first
    \draw[thick] (-2, \yone) -- (\ex, \ey-1);
    \fill (-2, \yone) circle [x radius=0.15, y radius=0.15];    
    \node at (-2, \yone+0.5) {$\mlinstance_{0}=1$};
    \node at (-1,{(\yone+0.4)/1.5}) {\figthetazero};

    % second
    \draw[thick] (-2, \ytwo) -- (\ex, \ey-1);
    \fill (-2, \ytwo) circle [x radius=0.15, y radius=0.15];    
    \node at (-2, \ytwo+0.5) {$\mlinstance_{1}$};
    \node at (-1,{(\ytwo+0.4)/1.5}) {\figthetaone};

    % third
    \draw[thick] (-2, \ythree) -- (\ex, \ey-1);
    \fill (-2, \ythree) circle [x radius=0.15, y radius=0.15];    
    \node at (-2, \ythree+0.5) {$\mlinstance_{2}$};
    \node at (-1,{(\ythree+0.4)/2}) {\figthetatwo};
    
    \draw[fill=PastelBlue] (\ex,\ey) arc(90:270:\eradius) --cycle;
    \draw[fill=PastelBlue2] (\ex,\ey) arc(90:-90:\eradius) --cycle;
    
    \node[font=\Large] at (\ex - 0.4, 0) {\figpreactivation}; %
    \node[font=\Large] at (\ex + 0.5, 0) {\figactivation};

    \draw[thick] (\ex+1, 0) -- (\ex+2, 0);
\end{tikzpicture}
    }    
    \end{column}
    \begin{column}{.4\textwidth}
    \centering
    \providecommand{\figfuncdecision}{x*x}

    \resizebox{\linewidth}{!}{
        \begin{tikzpicture}
\providecommand{\figfuncdecision}{-x+1}
\providecommand{\figtopcolor}{PastelRed}
\providecommand{\figbottomcolor}{PastelBlue}
\providecommand{\figvertical}{}
\providecommand{\figleftcolor}{PastelBlue}
\providecommand{\figrightcolor}{PastelRed}
\def\figemptytest{}
\begin{axis}[
    legend pos=south east,% <-- this line controls the width
    xlabel={$\mlinstance_1$},
    ylabel={$\mlinstance_2$},
    xmin=-1.5, xmax=1.5,
    ymin=-1.5, ymax=1.5,
    xtick={-1, 0, 1},
    ytick={-1, 0, 1},
    ymajorgrids=true,
    xmajorgrids=true,
    grid style=dashed,
]
    \ifx\figvertical\figemptytest
    \addplot [
        name path=f,
        domain=-1.5:1.5,
        samples=2,
        color=black,
        dashed,
        thick,
        ]
        {\figfuncdecision};
    \path [name path=bottom](\pgfkeysvalueof{/pgfplots/xmin},\pgfkeysvalueof{/pgfplots/ymin}) --
    (\pgfkeysvalueof{/pgfplots/xmax},\pgfkeysvalueof{/pgfplots/ymin});

    \path [name path=top]
    (\pgfkeysvalueof{/pgfplots/xmin},\pgfkeysvalueof{/pgfplots/ymax}) --
    (\pgfkeysvalueof{/pgfplots/xmax},\pgfkeysvalueof{/pgfplots/ymax});


    \addplot[\figtopcolor, opacity=0.3] fill between [of=f and top];
    \addplot[\figbottomcolor, opacity=0.3] fill between [of=f and bottom];
    \else
    \fill[\figleftcolor, opacity=0.3] (axis cs:-1.5,-1.5) rectangle (axis cs:\figvertical,1.5);
    \fill[\figrightcolor, opacity=0.3] (axis cs:\figvertical,-1.5) rectangle (axis cs:1.5,1.5);
    \draw[black, dashed, thick] (axis cs:\figvertical,-1.5) -- (axis cs:\figvertical,1.5);
    \fi

    \addplot[
        only marks,
        ultra thick,
         color=red,
         mark=-,
         mark options={scale=3}
        ]
        coordinates {
        (-1,-1)
        };
    \addplot[
        only marks,
        ultra thick,
         color=blue,
         mark=+,
         mark options={scale=3}
        ]
        coordinates {
        (1,1)(-1,1)(1,-1)
        };
\end{axis}
\end{tikzpicture}

    }
    \end{column}
\end{columns}
    
\end{frame}

\begin{frame}{Algoritmo de Aprendizado do Perceptron (PLA)}
% Inicialize


% %first data
% \tikzarrow{10.75}{5.2}
% \tikzarrow{13.75}{5.2}
% \tikzarrow{10.75}{2.6}
% \tikzarrow{13.75}{2.6}
% \tikzarrow{10}{-1.7}
% \tikzarrow{13}{-1.7}
% \tikzarrow{13}{-4.2}
% \tikzarrow{10}{-4.2}
    % \only<1>{
    % }
    % \only<2,5,8,11,14>{
    % }
    % \only<3,6,9,12,15>{
    % }
    % \only<4,7,10,13,16>{
    % }
    
%


    % \only<3,4,5,6>{
    % }
    % \only<7,8,9>{
    %     \tikzarrow{13}{-1.7}
    % }
    % \only<10,11,12>{
    %     \tikzarrow{13}{-4.2}
    %     }
    % \only<13,14,15>{
    %     \tikzarrow{10}{-4.2}
    % }

    
%

\begin{columns}
    \begin{column}{.6\textwidth}
    \visible<1>{
        \tikzarrow{0.25}{7.1}
    }
    \visible<2,5,8,11,14>{
        \tikzarrow{0.25}{6.65}
    }
    \visible<3,6,9,12>{
        \tikzarrow{0.25}{6.15}
    }
    \visible<4,7,10,13>{
        \tikzarrow{0.25}{5.7}
    }
    \visible<15>{
        \tikzarrow{0.25}{5.15}
    }
    
        \begin{algorithmic}[1]
        \State Inicialize os pesos $\thetav$ aleatoriamente
        \While{existirem observações mal classificadas}
          \State Selecione $\X^{(i)}$ tal que $\text{sinal}(\X^{(i)} \thetav^\top) \neq \y^{(i)}$
          \State Atualize: $\thetav \gets \thetav + \y^{(i)} \X^{(i)}$
        \EndWhile
        \end{algorithmic}
        \vspace{10pt}
        
        \only<1-3>{
            \providecommand{\figthetazero}{$-1$}
            \providecommand{\figthetaone}{$-1$}
            \providecommand{\figthetatwo}{$-1$}
            \(\thetav^\top=[-1,-1,-1]\)
        }
        \only<4>{
            \providecommand{\figthetazero}{$0$}
            \providecommand{\figthetaone}{$-2$}
            \providecommand{\figthetatwo}{$0$}
            \(\thetav^\top=[0,-2,0]=[-1,-1,-1]+[1,-1,1]\)
        }
        \only<5-6>{
            \providecommand{\figthetazero}{$0$}
            \providecommand{\figthetaone}{$-2$}
            \providecommand{\figthetatwo}{$0$}
            \(\thetav^\top=[0,-2,0]\)
        }
        \only<7>{
            \providecommand{\figthetazero}{$1$}
            \providecommand{\figthetaone}{$-1$}
            \providecommand{\figthetatwo}{$1$}
            \(\thetav^\top=[1,-1,1]=[0,-2,0]+[1,1,1]\)
        }
        \only<8-9>{
            \providecommand{\figthetazero}{$1$}
            \providecommand{\figthetaone}{$-1$}
            \providecommand{\figthetatwo}{$1$}
            \(\thetav^\top=[1,-1,1]\)
        }
        \only<10>{
            \providecommand{\figthetazero}{$2$}
            \providecommand{\figthetaone}{$0$}
            \providecommand{\figthetatwo}{$0$}
            \(\thetav^\top=[2,0,0]=[1,-1,1]+[1,1,-1]\)
        }
        \only<11-12>{
            \providecommand{\figthetazero}{$2$}
            \providecommand{\figthetaone}{$0$}
            \providecommand{\figthetatwo}{$0$}
            \(\thetav^\top=[2,0,0]\)
        }
        \only<13>{
            \providecommand{\figthetazero}{$1$}
            \providecommand{\figthetaone}{$1$}
            \providecommand{\figthetatwo}{$1$}
            \(\thetav^\top=[1,1,1]=[2,0,0]+(-1)[1,-1,-1]\)
        }
        \only<14->{
            \providecommand{\figthetazero}{$1$}
            \providecommand{\figthetaone}{$1$}
            \providecommand{\figthetatwo}{$1$}
            \(\thetav^\top=[1,1,1]\)
        }
        \\
    \centering
    \resizebox{.5\linewidth}{!}{
        \begin{tikzpicture}
    \providecommand{\figpreactivation}{$z^{(i)}$}
    \providecommand{\figactivation}{$a^{(i)}$}
    \providecommand{\figthetazero}{$\bm{\theta_{0}}$}
    \providecommand{\figthetaone}{$\bm{\theta_{1}}$}
    \providecommand{\figthetatwo}{$\bm{\theta_{2}}$}
    \pgfmathsetmacro{\ex}{0.5}
    \pgfmathsetmacro{\ey}{1}
    \pgfmathsetmacro{\eradius}{1}
    \pgfmathsetmacro{\yone}{1.5}
    \pgfmathsetmacro{\ytwo}{0}
    \pgfmathsetmacro{\ythree}{-1.5}
    % \draw[help lines] (-2,-2) grid (3,3);
    % first
    \draw[thick] (-2, \yone) -- (\ex, \ey-1);
    \fill (-2, \yone) circle [x radius=0.15, y radius=0.15];    
    \node at (-2, \yone+0.5) {$\mlinstance_{0}=1$};
    \node at (-1,{(\yone+0.4)/1.5}) {\figthetazero};

    % second
    \draw[thick] (-2, \ytwo) -- (\ex, \ey-1);
    \fill (-2, \ytwo) circle [x radius=0.15, y radius=0.15];    
    \node at (-2, \ytwo+0.5) {$\mlinstance_{1}$};
    \node at (-1,{(\ytwo+0.4)/1.5}) {\figthetaone};

    % third
    \draw[thick] (-2, \ythree) -- (\ex, \ey-1);
    \fill (-2, \ythree) circle [x radius=0.15, y radius=0.15];    
    \node at (-2, \ythree+0.5) {$\mlinstance_{2}$};
    \node at (-1,{(\ythree+0.4)/2}) {\figthetatwo};
    
    \draw[fill=PastelBlue] (\ex,\ey) arc(90:270:\eradius) --cycle;
    \draw[fill=PastelBlue2] (\ex,\ey) arc(90:-90:\eradius) --cycle;
    
    \node[font=\Large] at (\ex - 0.4, 0) {\figpreactivation}; %
    \node[font=\Large] at (\ex + 0.5, 0) {\figactivation};

    \draw[thick] (\ex+1, 0) -- (\ex+2, 0);
\end{tikzpicture}
    }    
    \end{column}
    \begin{column}{.4\textwidth}
    \only<3,4>{
        \tikzarrow{10.75}{5.9}
    }
    \only<6,7>{
        \tikzarrow{13.75}{5.9}
    }
    \only<9,10>{
        \tikzarrow{13.75}{3.4}
    }
    \only<12,13>{
        \tikzarrow{10.75}{3.4}
    }
    \centering
    \only<1,2,3>{
        \providecommand{\figfuncdecision}{-x-1}
    }
    \only<4,5,6>{
        \providecommand{\figvertical}{0}
        \providecommand{\figleftcolor}{PastelBlue}
        \providecommand{\figrightcolor}{PastelRed}
    }
    \only<7,8,9>{
        \providecommand{\figtopcolor}{PastelBlue}
        \providecommand{\figbottomcolor}{PastelRed}
        \providecommand{\figfuncdecision}{x-1}

    }
    \only<10,11,12>{
        \providecommand{\figtopcolor}{PastelRed}
        \providecommand{\figbottomcolor}{PastelBlue}
        \providecommand{\figfuncdecision}{x*x}
    }
    \only<13->{
        \providecommand{\figtopcolor}{PastelBlue}
        \providecommand{\figbottomcolor}{PastelRed}
        \providecommand{\figfuncdecision}{-x-1}
    }

    \resizebox{\linewidth}{!}{
        \begin{tikzpicture}
\providecommand{\figfuncdecision}{-x+1}
\providecommand{\figtopcolor}{PastelRed}
\providecommand{\figbottomcolor}{PastelBlue}
\providecommand{\figvertical}{}
\providecommand{\figleftcolor}{PastelBlue}
\providecommand{\figrightcolor}{PastelRed}
\def\figemptytest{}
\begin{axis}[
    legend pos=south east,% <-- this line controls the width
    xlabel={$\mlinstance_1$},
    ylabel={$\mlinstance_2$},
    xmin=-1.5, xmax=1.5,
    ymin=-1.5, ymax=1.5,
    xtick={-1, 0, 1},
    ytick={-1, 0, 1},
    ymajorgrids=true,
    xmajorgrids=true,
    grid style=dashed,
]
    \ifx\figvertical\figemptytest
    \addplot [
        name path=f,
        domain=-1.5:1.5,
        samples=2,
        color=black,
        dashed,
        thick,
        ]
        {\figfuncdecision};
    \path [name path=bottom](\pgfkeysvalueof{/pgfplots/xmin},\pgfkeysvalueof{/pgfplots/ymin}) --
    (\pgfkeysvalueof{/pgfplots/xmax},\pgfkeysvalueof{/pgfplots/ymin});

    \path [name path=top]
    (\pgfkeysvalueof{/pgfplots/xmin},\pgfkeysvalueof{/pgfplots/ymax}) --
    (\pgfkeysvalueof{/pgfplots/xmax},\pgfkeysvalueof{/pgfplots/ymax});


    \addplot[\figtopcolor, opacity=0.3] fill between [of=f and top];
    \addplot[\figbottomcolor, opacity=0.3] fill between [of=f and bottom];
    \else
    \fill[\figleftcolor, opacity=0.3] (axis cs:-1.5,-1.5) rectangle (axis cs:\figvertical,1.5);
    \fill[\figrightcolor, opacity=0.3] (axis cs:\figvertical,-1.5) rectangle (axis cs:1.5,1.5);
    \draw[black, dashed, thick] (axis cs:\figvertical,-1.5) -- (axis cs:\figvertical,1.5);
    \fi

    \addplot[
        only marks,
        ultra thick,
         color=red,
         mark=-,
         mark options={scale=3}
        ]
        coordinates {
        (-1,-1)
        };
    \addplot[
        only marks,
        ultra thick,
         color=blue,
         mark=+,
         mark options={scale=3}
        ]
        coordinates {
        (1,1)(-1,1)(1,-1)
        };
\end{axis}
\end{tikzpicture}

    }
    \end{column}
\end{columns}
    
\end{frame}






\begin{frame}{Considerações sobre o PLA}
\begin{columns}
    \begin{column}{.5\linewidth}
        \begin{itemize}
          \item Descoberto empiricamente
          \item Intuição geométrica: produto escalar
          \begin{itemize}
              \item positivo quando $\alpha<90^{\circ}$
              \item negativo quando $\alpha>90^{\circ}$
          \end{itemize}
          \item A regra de atualização de $\thetav$ ajusta o ângulo do vetor de pesos em relação às instâncias
          \item Sempre converge para problemas linearmente separáveis
        \end{itemize}
    \end{column}
    \begin{column}{.5\linewidth}
\begin{tikzpicture}[scale=2, >=Latex]
    % Axes
    \draw[->, gray!60] (-1.5,0) -- (1.5,0) node[right] {};
    \draw[->, gray!60] (0,-.2) -- (0,1.5) node[above] {};

    % Vectors u and v
    \coordinate (O) at (0,0);
    \coordinate (U) at (1.5,0.5);
    \coordinate (V) at (0.5,1.2);
    \coordinate (X) at (-1,0.7);

    \draw[->, thick, blue] (O) -- (U) node[above right] {$\X^{(i)}$};
    \draw[->, thick, red] (O) -- (V) node[above left] {$\thetav$};
    \draw[->, thick, orange] (O) -- (X) node[above left] {$\thetav-\X^{(i)}$};
    \draw[-, dotted, thick, gray!60] (V) -- (X) node[above left] {};
    \pic [draw=green!50!black, fill=green!20, angle eccentricity=1.3, angle radius=1cm] {angle = U--O--V};
    \node at (0.2,0.2) {$\alpha$};
    \node at (-0.4,1.2) {$\yhat=-1$};
    % Angle between vectors
\end{tikzpicture}

    \end{column}
\end{columns}
\end{frame}

\section{O problema XOR}
\begin{frame}{Limitação do Perceptron}
\begin{itemize}
  \item O Perceptron é um modelo linear (traça um hiperplano como fronteira de decisão)
  \item Marvin Minsky e Seymour Papert salientaram essa limitação com o problema  XOR\footfullcite{minsky1969perceptron}
\end{itemize}
\vspace{5pt}
\centering
\providecommand{\figtopcolor}{white}
\providecommand{\figbottomcolor}{white}
\providecommand{\figfuncdecision}{x*x}
\resizebox{.38\linewidth}{!}{
\input{slides/img/tikz/perceptronfronteiradecisaoxor.tikz}
}

\end{frame}




\begin{frame}{Vamos observar esses dois Perceptrons}
\begin{columns}
  \column{0.5\textwidth}
  \centering
  \providecommand{\figthetazero}{$0.5$}
  \providecommand{\figthetaone}{$1$}
  \providecommand{\figthetatwo}{$1$}
  \resizebox{.6\linewidth}{!}{
  \begin{tikzpicture}
    \providecommand{\figpreactivation}{$z^{(i)}$}
    \providecommand{\figactivation}{$a^{(i)}$}
    \providecommand{\figthetazero}{$\bm{\theta_{0}}$}
    \providecommand{\figthetaone}{$\bm{\theta_{1}}$}
    \providecommand{\figthetatwo}{$\bm{\theta_{2}}$}
    \pgfmathsetmacro{\ex}{0.5}
    \pgfmathsetmacro{\ey}{1}
    \pgfmathsetmacro{\eradius}{1}
    \pgfmathsetmacro{\yone}{1.5}
    \pgfmathsetmacro{\ytwo}{0}
    \pgfmathsetmacro{\ythree}{-1.5}
    % \draw[help lines] (-2,-2) grid (3,3);
    % first
    \draw[thick] (-2, \yone) -- (\ex, \ey-1);
    \fill (-2, \yone) circle [x radius=0.15, y radius=0.15];    
    \node at (-2, \yone+0.5) {$\mlinstance_{0}=1$};
    \node at (-1,{(\yone+0.4)/1.5}) {\figthetazero};

    % second
    \draw[thick] (-2, \ytwo) -- (\ex, \ey-1);
    \fill (-2, \ytwo) circle [x radius=0.15, y radius=0.15];    
    \node at (-2, \ytwo+0.5) {$\mlinstance_{1}$};
    \node at (-1,{(\ytwo+0.4)/1.5}) {\figthetaone};

    % third
    \draw[thick] (-2, \ythree) -- (\ex, \ey-1);
    \fill (-2, \ythree) circle [x radius=0.15, y radius=0.15];    
    \node at (-2, \ythree+0.5) {$\mlinstance_{2}$};
    \node at (-1,{(\ythree+0.4)/2}) {\figthetatwo};
    
    \draw[fill=PastelBlue] (\ex,\ey) arc(90:270:\eradius) --cycle;
    \draw[fill=PastelBlue2] (\ex,\ey) arc(90:-90:\eradius) --cycle;
    
    \node[font=\Large] at (\ex - 0.4, 0) {\figpreactivation}; %
    \node[font=\Large] at (\ex + 0.5, 0) {\figactivation};

    \draw[thick] (\ex+1, 0) -- (\ex+2, 0);
\end{tikzpicture}
  }
  
  \visible<2>{\textbf{OR}}\\
  \begin{tabular}{|r|r|r|}
  \hline
  $\X_1$ & $\X_2$ & $\y$ \\
  \hline
  -1 & -1 & \visible<2>{\textcolor{red}{-1}} \\
  \hline
  -1 & 1  & \visible<2>{\textcolor{blue}{1}} \\
  \hline
   1 & -1 & \visible<2>{\textcolor{blue}{1}} \\
   \hline
   1 & 1  & \visible<2>{\textcolor{blue}{1}} \\
   \hline
  \end{tabular}

  \column{0.5\textwidth}
  \centering
  \providecommand{\figthetazero}{$-1.5$}
  \providecommand{\figthetaone}{$1$} \providecommand{\figthetatwo}{$1$}
  \resizebox{.6\linewidth}{!}{
  \begin{tikzpicture}
    \providecommand{\figpreactivation}{$z^{(i)}$}
    \providecommand{\figactivation}{$a^{(i)}$}
    \providecommand{\figthetazero}{$\bm{\theta_{0}}$}
    \providecommand{\figthetaone}{$\bm{\theta_{1}}$}
    \providecommand{\figthetatwo}{$\bm{\theta_{2}}$}
    \pgfmathsetmacro{\ex}{0.5}
    \pgfmathsetmacro{\ey}{1}
    \pgfmathsetmacro{\eradius}{1}
    \pgfmathsetmacro{\yone}{1.5}
    \pgfmathsetmacro{\ytwo}{0}
    \pgfmathsetmacro{\ythree}{-1.5}
    % \draw[help lines] (-2,-2) grid (3,3);
    % first
    \draw[thick] (-2, \yone) -- (\ex, \ey-1);
    \fill (-2, \yone) circle [x radius=0.15, y radius=0.15];    
    \node at (-2, \yone+0.5) {$\mlinstance_{0}=1$};
    \node at (-1,{(\yone+0.4)/1.5}) {\figthetazero};

    % second
    \draw[thick] (-2, \ytwo) -- (\ex, \ey-1);
    \fill (-2, \ytwo) circle [x radius=0.15, y radius=0.15];    
    \node at (-2, \ytwo+0.5) {$\mlinstance_{1}$};
    \node at (-1,{(\ytwo+0.4)/1.5}) {\figthetaone};

    % third
    \draw[thick] (-2, \ythree) -- (\ex, \ey-1);
    \fill (-2, \ythree) circle [x radius=0.15, y radius=0.15];    
    \node at (-2, \ythree+0.5) {$\mlinstance_{2}$};
    \node at (-1,{(\ythree+0.4)/2}) {\figthetatwo};
    
    \draw[fill=PastelBlue] (\ex,\ey) arc(90:270:\eradius) --cycle;
    \draw[fill=PastelBlue2] (\ex,\ey) arc(90:-90:\eradius) --cycle;
    
    \node[font=\Large] at (\ex - 0.4, 0) {\figpreactivation}; %
    \node[font=\Large] at (\ex + 0.5, 0) {\figactivation};

    \draw[thick] (\ex+1, 0) -- (\ex+2, 0);
\end{tikzpicture}
  }
  
  \visible<2>{\textbf{AND:}}\\
  \begin{tabular}{|r|r|r|}
  \hline
  $\X_1$ & $\X_2$ & $\y$ \\
  \hline
  -1 & -1 & \visible<2>{\textcolor{red}{-1}} \\
  \hline
  -1 & 1  & \visible<2>{\textcolor{red}{-1}} \\
  \hline
   1 & -1 & \visible<2>{\textcolor{red}{-1}} \\
   \hline
   1 & 1  & \visible<2>{\textcolor{blue}{1}} \\
   \hline
  \end{tabular}
\end{columns}
\end{frame}

\begin{frame}{Resolvendo XOR}
\begin{columns}
    \begin{column}{0.4\linewidth}
        \begin{tikzpicture}
            \pgfmathsetmacro{\ex}{0.5}
            \pgfmathsetmacro{\ey}{1.}
            \pgfmathsetmacro{\eradius}{0.5}
            \pgfmathsetmacro{\yone}{1.5}
            \pgfmathsetmacro{\ytwo}{0}
            \pgfmathsetmacro{\ythree}{-1.5}
            % \draw[help lines] (-2,-2) grid (3,3);
            % first
            \fill (-2, \yone) circle [x radius=0.15, y radius=0.15];    
            \node at (-2, \yone+0.5) {$\mlinstance_{0}=1$};
        
            % second
            \fill (-2, \ytwo) circle [x radius=0.15, y radius=0.15];    
            \node at (-2, \ytwo+0.5) {$\mlinstance_{1}$};
        
            % third
            \fill (-2, \ythree) circle [x radius=0.15, y radius=0.15];    
            \node at (-2, \ythree+0.5) {$\mlinstance_{2}$};
            


        \visible<1,3>{
            \draw[thick] (-2, \yone) -- (\ex, \ey-1);
            \draw[thick] (-2, \ytwo) -- (\ex, \ey-1);
            \draw[thick] (-2, \ythree) -- (\ex, \ey-1);
            \node at (-1.5,{\yone}) {\small \textcolor{blue}{$-1.5$}};
            \node at (-1.5,{\ytwo+0.2}) {\small \textcolor{blue}{$1$}};
            \node at (-1.5,{\ythree+0.6}) {\small \textcolor{blue}{$-1$}};
            \draw[thick] (\ex, 0) -- (\ex+2, 0);
            \filldraw[color=black,fill=PastelBlue](\ex,\ey-1) circle (0.5);
            \node[font=\tiny] at (\ex, 0) {$\varphi(z^{(i)})$};
            
            }           

        \visible<2->{
            \node at (-1.5,{\yone-1}) {\small \textcolor{green}{$-1.5$}};
            \node at (-1.5,{\ytwo-0.2}) {\small \textcolor{green}{$-1$}};
            \node at (-1.5,{\ythree-0.3}) {\small \textcolor{green}{$1$}};
            \draw[thick] (-2, \yone) -- (\ex, \ythree);
            \draw[thick] (-2, \ytwo) -- (\ex, \ythree);
            \draw[thick] (-2, \ythree) -- (\ex, \ythree);
        
            \draw[thick] (\ex,\ythree) -- (\ex+2, 0);
            
            \filldraw[color=black,fill=PastelGreen](\ex,\ythree) circle (0.5);

            \node[font=\tiny] at (\ex, \ythree) {$\varphi(z^{(i)})$};%
            }
            \visible<3>{
                \draw[thick] (-2, \yone) -- (\ex+2, 0);
                \node at (1.7,{\yone-1}) {\small \textcolor{red}{$0.5$}};
                \node at (1.7,{\ytwo-0.2}) {\small \textcolor{red}{$1$}};
                \node at (1.7,{\ytwo-0.8}) {\small \textcolor{red}{$1$}};
                
                \draw[thick,->] (\ex+2, 0) -- (\ex+3, 0);
                \filldraw[color=black,fill=PastelRed]({\ex+2},{\ey-1}) circle (0.5);
                \node[font=\tiny] at (\ex+2, 0) {$\varphi(z^{(i)})$};%
            }
        \end{tikzpicture}
    \end{column}
    \begin{column}{0.6\linewidth}
    \centering
    \resizebox{.9\linewidth}{!}{
    \begin{tikzpicture}
        \begin{axis}[
            legend pos=south east,% <-- this line controls the width
            xlabel={$\mlinstance_1$},
            ylabel={$\mlinstance_2$},
            xmin=-1.5, xmax=1.5,
            ymin=-1.5, ymax=1.5,
            xtick={-1, 0, 1},
            ytick={-1, 0, 1},
            ymajorgrids=true,
            xmajorgrids=true,
            grid style=dashed,
        ]
            
                
            \path [name path=bottom](\pgfkeysvalueof{/pgfplots/xmin},\pgfkeysvalueof{/pgfplots/ymin}) --
            (\pgfkeysvalueof{/pgfplots/xmax},\pgfkeysvalueof{/pgfplots/ymin});
            
            \path [name path=top]
            (\pgfkeysvalueof{/pgfplots/xmin},\pgfkeysvalueof{/pgfplots/ymax}) --
            (\pgfkeysvalueof{/pgfplots/xmax},\pgfkeysvalueof{/pgfplots/ymax});
            
            \only<2>{
            \addplot [
                name path=g,
                domain=-1.5:1.5, 
                samples=2, 
                color=black,
                dashed,
                thick,
                ]
                {x+1};    
                \addplot[PastelBlue, opacity=0.3] fill between [of=g and top]; 
                \addplot[PastelRed, opacity=0.3] fill between [of=g and bottom]; 
            }
            \only<1>{
            \addplot [
                    name path=f,
                    domain=-1.5:1.5, 
                    samples=2, 
                    color=black,
                    dashed,
                    thick,
                    ]
                    {x-1};
            
                \addplot[PastelRed, opacity=0.3] fill between [of=f and top]; 
                \addplot[PastelBlue, opacity=0.3] fill between [of=f and bottom]; 
            }
            \only<3>{
            \addplot [
                    name path=f,
                    domain=-1.5:1.5, 
                    samples=2, 
                    color=black,
                    dashed,
                    thick,
                    ]
                    {x-1};
            \addplot [
                name path=g,
                domain=-1.5:1.5, 
                samples=2, 
                color=black,
                dashed,
                thick,
                ]
                {x+1};
                \addplot[PastelBlue, opacity=0.3] fill between [of=g and top]; 
                \addplot[PastelRed, opacity=0.3] fill between [of=f and g]; 
                \addplot[PastelBlue, opacity=0.3] fill between [of=f and bottom]; 
            }
            \addplot[
                only marks,
                ultra thick,
                 color=red,
                 mark=-,
                 mark options={scale=3}
                ]
                coordinates {
                (-1,-1)(1,1)
                };
            \addplot[
                only marks,
                ultra thick,
                 color=blue,
                 mark=+,
                 mark options={scale=3}
                ]
                coordinates {
                (-1,1)(1,-1)
                };
        \end{axis}
        \end{tikzpicture}}
    \end{column}
\end{columns}
    
\end{frame}

\begin{frame}{Resolvendo XOR}
\begin{itemize}
    \item Quando organizamos vários Perceptrons em camadas chamamos de \textit{MultiLayer Perceptron} (MLP)
    \item PLA não funciona para MLP
    \item Essa limitação minou o interesse por redes neurais durante anos (Inverno da IA)
    \item Para treinar um MLP vamos precisar usar retropropagação de erros (\textit{backpropagation}) e descida de gradiente (\textit{gradient descent})
\end{itemize}
    
\end{frame}
\section*{Referências e Leituras}
%------------------------------------------------
\begin{frame}{Leituras Recomendadas}
    \leiturasize
    \begin{itemize}
        \item Capítulo 1: \fullcite{abu-mostafa_learning_2012}
        \item Capítulo 1: \fullcite{haykin2009neural}
    \end{itemize}

\end{frame}



%------------------------------------------------

\begin{frame}{Referências}
    \printbibliography

\end{frame}

%------------------------------------------------

% \begin{frame}
%     \Huge{\centerline{\textbf{The End}}}
% \end{frame}

%----------------------------------------------------------------------------------------

\end{document}